% ===== 原书 p.39 =====
% 注：p.39 上半页属 §4 末例（圆的例子），按节界归 c06；本块自 \section 起。
\section{Bott 周期性的一个证明}
\subsection{预备}

正文始终用上同调、K 理论与表示论这套工具，从分析一侧求解指标问题。本附录把方向反过来：把指标理论用到拓扑上。具体做法是用一族椭圆算子的指标，给 Bott 周期性定理一个简单证明。§1 已说明，Bott 周期性定理是 K 理论的基本基石，使它具有周期上同调理论的结构。希望这个证明能让读者看清代数拓扑与椭圆算子之间的根本联系（若正文其余部分还没能让读者看清）。这一联系的更多材料见 [3]。我们的证明沿 [5] 的思路，是 [4] 中更一般论证的一个稍简情形。本附录后段将简述：如何把这一般论证推广，去证明 Thom 同构定理的等变版本。

开始前先回顾 §1 中 Bott 周期性定理的陈述。设 $b$ 是 $\tilde K(S^2)=K(\C)$ 的生成元。注意外张量积（exterior tensor product）$\BT:K(X)\otimes K(Y)\to K(X\times Y)$ 限制为乘积 $\BT:\tilde K(X)\otimes\tilde K(Y)\to\tilde K(X\wedge Y)$，\footnote{利用分解 $K(X)\iso\tilde K(X)\oplus K(\pt)$，外张量积分裂为
\[
\otimes:\bigl(\tilde K(X)\otimes\tilde K(Y)\bigr)\oplus\bigl(\tilde K(X)\oplus\tilde K(Y)\bigr)\oplus K(\pt)\longrightarrow\tilde K(X\wedge Y)\oplus\tilde K(X\vee Y)\oplus K(\pt).
\]
由 $\tilde K(X)\oplus\tilde K(Y)\to\tilde K(X\vee Y)$ 可知，外张量积限制为乘积映射
\[
\BT:\tilde K(X)\otimes\tilde K(Y)\to\tilde K(S\wedge Y),
\]
其中回忆 $X\vee Y=(X\times\pt)\cup(\pt\times Y)$ 是单点并，$X\wedge Y=(X\times Y)/(X\vee Y)$ 是 smash 积。} 据此定义同态
\[
\beta:K(X)\to K^{-2}(X)=\tilde K(S^2\wedge X^+).
\]

\transnote{本附录原文有若干处符号笔误，译文一律照原样排印。已知的有：上面脚注中的 $\tilde K(S\wedge Y)$ 应为 $\tilde K(X\wedge Y)$；下一段开头的 $\beta(X)$ 应为 $\beta(x)$；p.40 的 $\sum_{0\le m<n}b_m e^{imx}$ 应为 $e^{im\theta}$；p.42 的 $f(x)=z^{-2}$ 应为 $f(z)$；p.44 的 $\alpha(X)\cdot y$ 应为 $\alpha(x)\cdot y$，同页 $u\cdot\lambda_1$ 应为 $u\cdot b$，$K(X\times\R^n)$ 应为 $K(X\times\R^2)$；同页 $\alpha\beta(X)$ 应为 $\alpha\beta(x)$。读到这些位置按上下文理解即可。}

% ===== 原书 p.40 =====
映射由 $\beta(X)=b\BT x$ 给出。Bott 周期性定理断言这个映射是同构。用紧支撑 K 理论注意到
\[
K(\R^n\times X)=\tilde K\bigl((\R^n\times X)^+\bigr)=\tilde K(S^n\wedge X^+)=K^{-n}(X),
\]
于是可以把 $\beta$ 看成同态
\[
\beta:K(X)\xrightarrow{\,b\BT\cdot\,}K(\R^2\times X),
\]
其中 $b$ 视为 $K(\R^2)$ 的元素。为证明 Bott 周期性定理，我们将用一族椭圆算子的指标，构造逆同态 $\alpha:K(\R^2\times X)\to K(X)$。

\subsection{Wiener--Hopf 算子}

先看圆周 $S^1=\R/2\pi\Z$ 上一个具体的伪微分算子。§3 引入过 $S^1$ 上复值函数的零阶伪微分算子 $P$，它在 Fourier 级数下由
\[
P(e^{ik\theta})=\begin{cases} e^{ik\theta} & k\ge 0,\\ 0 & k<0 \end{cases}
\]
给出。这个算子可以看成从 $S^1$ 上复值函数空间到「只有正 Fourier 系数」的子空间的投影。设 $f:S^1\to\C^*$ 连续，其中 $\C^*$ 是非零复数。可以构造相应的伪微分算子 $P_f=PfP+(1-P)$。写 $f(\theta)=\sum_n a_n e^{in\theta}$，按 Fourier 级数展开得到
\[
P_f(e^{ik\theta})=\begin{cases} \sum_{n>-k}a_n e^{i(n+k)\theta} & k\ge 0,\\ e^{ik\theta} & k<0, \end{cases}
\]
注意 $P_f$ 是椭圆的，因为它的零阶象征由
\[
\sigma_{P_f}(\theta,\xi)=\begin{cases} f(\theta) & \xi>0,\\ 1 & \xi<0 \end{cases}
\]
给出。

\transnote{$P$ 的像含常数项：上图对 $k=0$ 也取 $e^{ik\theta}$，故原文「只有正 Fourier 系数」应读作「非负」。这也暴露下一式的下界问题：按 $P_f=PfP+(1-P)$ 逐项计算，求和应为 $\sum_{n\ge-k}$；原文写 $\sum_{n>-k}$，在 $k=0$ 时会漏掉 $a_0$，而原文随后断言 $\Coker P_n$ 为 $n$ 维，只有用 $\sum_{n\ge-k}$ 才对得上。此处按上下文补边界项理解。}

特别地，令 $P_n=P_f$，其中 $f(\theta)=e^{in\theta}$。当 $n\ge 0$ 时有
\[
P_n(e^{ik\theta})=\begin{cases} e^{i(n+k)\theta} & k\ge 0,\\ e^{ik\theta} & k<0, \end{cases}
\]
于是 $\Ker P_n$ 平凡，而 $\Coker P_n$ 是 $n$ 维空间，由形如 $g(\theta)=\sum_{0\le m<n}b_m e^{imx}$ 的函数组成。故 $\Ind P_n=-n$。同理，当 $n<0$ 时有
\[
P_n(e^{ik\theta})=\begin{cases} e^{i(n+k)\theta} & k\ge -n,\\ 0 & 0\le k<-n,\\ e^{ik\theta} & k<0. \end{cases}
\]
此时 $\Coker P_n$ 平凡，而 $\Ker P_n\iso\Coker P_{-n}$，于是同样得到 $\Ind P_n=-n$。

现在回顾：任何连续映射 $f:S^1\to\C^*$ 都同伦于形如 $f(\theta)=e^{in\theta}$ 的映射，映射的度定义为 $\deg f=n$。% ===== 原书 p.41 =====
指标是局部常值函数，于是 $\Ind P_f=-\deg f$。

现在推广这一构造。第一步，不妨把讨论扩充到 $S^1$ 上复向量值函数。设 $f:S^1\to\GL(N,\C)$ 连续，可以构造相应的椭圆算子 $P_f$。注意 $\Ind P_f$ 只依赖 $f$ 的同伦类；这里同伦类就是 $\GL(N,\C)$ 基本群中的类 $[f]\in\pi_1(\GL(N,\C))$。再进一步：设 $V$ 是紧空间 $X$ 上的复向量丛，并对每个 $x\in X$ 给出一族连续变化的连续映射 $f_x:S^1\to\Aut V_x$。于是得到 $S^1$ 上一族由 $X$ 参数化的伪微分算子 $P_f$，其中每个 $(P_f)_x=P_{f_x}$ 都椭圆。这样的集合称为椭圆族（elliptic family）。若 $\dim\Ker(P_f)_x$ 与 $\dim\Coker(P_f)_x$ 对所有 $x\in X$ 都是常数，则可把 $\Ker P_f$ 与 $\Coker P_f$ 定义为 $X$ 上相应的向量丛，从而定义
\[
\Ind P_f=[\Ker P_f]-[\Coker P_f]\in K(X).
\]
当 $\dim\Ker(P_f)_x$ 与 $\dim\Coker(P_f)_x$ 不是常数时，Fredholm 算子族的指标的一般构造见 [5, §2]。注意这个指标只依赖映射 $f:(x,\theta)\mapsto\Aut V_x$ 的同伦类 $[f]$。

\subsection{粘合构造}

设 $V$ 是 $X$ 上的向量丛，$f:(x,\theta)\mapsto\Aut V_x$ 是连续映射。如下构造 $S^2\times X$ 上的向量丛 $E$。设 $B^+$ 与 $B^-$ 是球面 $S^2$ 的闭上半球与闭下半球。二者之并 $B^+\cup B^-=S^2$ 是整个球面，交 $B^+\cap B^-=S^1$ 是赤道。取沿映射 $f$ 的并
\[
E=(B^-\times V)\cup_f(B^+\times V),
\]
其中沿交 $(B^-\times V)\cap(B^+\times V)=S^1\times V$ 作等同
\[
(\theta,v)\in B^-\times V_x\ \sim\ \bigl(\theta,f(\theta,x)v\bigr)\in B^+\times V_x.
\]
这个构造称为粘合构造（clutching construction），映射 $f:(x,\theta)\mapsto\Aut V_x$ 称为粘合函数（clutching function）。粘合函数的完整讨论见 [6, §1.4]。

注意 $S^2\times X$ 上每个向量丛 $E$ 都可由这一构造从 $X$ 上某个向量丛 $V$ 与适当的粘合函数 $f:(x,\theta)\mapsto\Aut V_x$ 得到\footnote{设 $i_+,i_-:X\to S^2\times X$ 是 $S^2$ 两个极点处的包含映射。注意 $B^+$ 与 $B^-$ 可缩，故 $E|_{B^+\times X}\iso B^+\times i_+^*E$，且 $E|_{B^-\times X}\iso B^+\times i_-^*E$。取 $V\iso i_-^*E\iso i_+^*E$，再把 $f:(x,\theta)\mapsto\Aut V_x$ 取为 $E|_{B^+\times X}$ 与 $E|_{B^-\times X}$ 在 $S^1\times X$ 上的等同。}；并且 $E$ 由 $f$ 的同伦类 $[f]$ 唯一决定（同构意义下，见 [6, §1.4.6]）。于是得到映射 $\Vect(S^2\times X)\to K(X)$：
\[
E\mapsto (V,[f])\mapsto P_f\mapsto\Ind P_f\in K(X),
\]
它再扩充为群同态 $\alpha:K(S^2\times X)\to K(X)$。注意\footnote{考虑空间序列 $X^+\vee Y^+\to X^+\times Y^+\to X^+\wedge Y^+=(X\times Y)^+$，得到正合列 $0\to K(X\times Y)\to\tilde K(X^+\times Y^+)\to K(X)\oplus K(Y)\to0$。特别地，取 $X=\R^n$、$Y=X$，有正合列 $0\to K(\R^n\times X)\to\tilde K(S^n\times X^+)\to\tilde K(S^n)\oplus K(X)\to0$。注意 $\tilde K(S^n\times X^+)=\tilde K(S^n)\oplus K(S^n\times X)$，可见 $K(S^n\times X)\iso K(\R^n\times X)\oplus K(X)$。} $K(S^2\times X)\iso K(\R^2\times X)\oplus K(X)$，取 $\alpha$ 在 $K(\R^2\times X)$ 上的限制，得到
\[
\alpha:K(\R^2\times X)\to K(X).
\]

\transnote{为什么先落到 $K(S^2\times X)$ 再取出 $K(\R^2\times X)$：粘合构造直接给出的丛住在 $S^2\times X$ 上，而 $K(S^2\times X)\iso K(\R^2\times X)\oplus K(X)$ 的两块按沿 $S^2$ 方向是否平凡分开——$K(X)$ 一块由 $X$ 上的丛沿 $S^2$ 平凡扩张而来，$K(\R^2\times X)$ 一块才携带沿 $S^2$ 的粘合信息。把 $\alpha$ 限制到后者，等于只看在 $S^2$ 方向非平凡的丛。Bott 映射的逆正是在这个直和项上构造的。}

% ===== 原书 p.42 =====
取 $X$ 为一点，则 $X$ 上的向量丛就是向量空间。对一维丛，粘合函数 $f:(x,\theta)\mapsto\Aut V_x$ 就是连续映射 $f:S^1\to\C^*$。于是对每个 $n\in\Z$，用 $V=\C$、$f(\theta)=e^{in\theta}$ 作粘合构造，得到 $S^2$ 上的丛 $E_n$。注意由构造立刻得到 $\alpha[E_n]=-n$。回顾每个这样的映射 $f$ 都同伦于形如 $f(\theta)=e^{in\theta}$ 的映射，可见这一构造给出 $S^2$ 上所有一维丛。特别地，这说明 $\tilde K(S^2)\iso\Z$。下面从这一观点看 $S^2$ 上向量丛的几个例子。

\begin{example}
把 $S^2\iso\C\cup\infty$ 看成 Riemann 球面，考虑复切丛 $T_\C S^2$。把 $B^-$ 等同于单位圆盘 $\{z\mid |z|\le1\}$，把 $B^+$ 等同于 $\{w\mid |w|\le1\}$，其中 $w=z^{-1}$（处理复结构时，希望坐标转移是复解析的）。点 $z\in B^-$ 处的切向量可写成 $\xi=d/dt\,f(t)|_{t=0}$，其中 $f(0)=z$。换到 $w$ 坐标，相应切向量为
\[
\xi_w=\frac{d}{dt}f(t)^{-1}\Big|_{t=0}=-f(t)^{-2}\frac{d}{dt}f(t)\Big|_{t=0}=-z^{-2}\xi_z.
\]
于是可以写
\[
T_\C S^2\iso\bigl((B^-\times\C)\cup(B^+\times\C)\bigr)\big/\sim,
\]
其中在单位圆 $\{z\mid |z|=1\}=B^-\cap B^+$ 上有等同
\[
(z,\xi)\in B^-\times\C\ \sim\ (z,-z^{-2}\xi)\in B^+\times\C.
\]
把单位圆转半圈可以看出，$f(z)=-z^{-2}$ 同伦于映射 $f(x)=z^{-2}$。这个映射对应 $f(\theta)=e^{-2i\theta}$，故 $T_\C S^2\iso E_{-2}$。
\end{example}

\begin{example}
考虑映射 $f(\theta)=e^{in\theta}$。相应椭圆伪微分算子 $P_n$ 的象征 $\sigma(P_n)\in K(TS^1)$ 由 $TS^1$ 上的复形
\[
0\longrightarrow TS^1\times\C\xrightarrow{\ \sigma_n\ }TS^1\times\C\longrightarrow 0
\]
给出，其中同态由
\[
\sigma_n(\theta,\xi)=\begin{cases} e^{in\theta} & \xi>0,\\ 1 & \xi<0 \end{cases}
\]
给出。设 $p^+\in B^+$、$p^-\in B^-$ 是球面 $S^2$ 的「北极」与「南极」。可把 $B^--\{p^-\}$ 等同于 $\{(\theta,\xi)\in TS^1\mid\xi\le0\}$，把 $B^+-\{p^+\}$ 等同于 $\{(\theta,\xi)\in TS^1\mid\xi\ge0\}$。注意 $\sigma_n(\theta,\xi)$ 同伦于同态 $\sigma'_n(\theta,\xi)$，后者由
\[
\sigma'(\theta,\xi)=\begin{cases} f(\theta)/(1+\xi) & \xi>0,\\ 1 & \xi<0 \end{cases}
\]
给出。可以把 $\sigma'_n$ 延拓为连续映射 $\sigma'_n:E_0\to E_n$，只需令
\[
\sigma'_n(p^-)=\lim_{\xi\to-\infty}\sigma'_n(\theta,\xi)=1,\qquad
\sigma'_n(p^+)=\lim_{\xi\to+\infty}\sigma'_n(\theta,\xi)=0.
\]
注意这个映射在圆 $S^1\times\{0\}$ 上连续：按粘合构造，我们把点 $((\theta,0),z)\in B^-\times\C$ 与 $((\theta,0),e^{in\theta}z)\in B^+\times\C$ 等同。于是象征复形延拓为 $S^2$ 上的复形
\[
0\longrightarrow E_0\xrightarrow{\ \sigma'_n\ }E_n\longrightarrow 0,
\]
并且 $\sigma(P_n)=[E_0]-[E_n]\in\tilde K(S^2)$。由此得到 $\alpha(\sigma(P_f))=n$。
\end{example}

% ===== 原书 p.43 =====
\begin{proposition}
对任意 $x\in K(TS^1)\iso K(\R^2)\iso\tilde K(S^2)$，有 $\aind_{S^1}(x)=-\alpha_\pt(x)$。
\end{proposition}

\begin{proof}
注意 $K(TS^1)\iso\Z$，故 $K(TS^1)$ 中每个元素 $x$ 都作为某个椭圆伪微分算子 $P_n$ 的象征 $x\in\sigma(P_n)$ 出现。于是由上面的讨论
\[
\begin{aligned}
\aind_{S^1}(x)&=\Ind P_n=-n\\
&=-\alpha_\pt\bigl(\sigma(P_n)\bigr)=-\alpha_\pt(x).
\end{aligned}
\]
\end{proof}

\begin{example}
回顾 §1：Bott 类 $\lambda_1\in K(\C)\iso\tilde K(S^2)$ 由 $\C$ 上的外复形 $\Lambda(\C)$ 表示：
\[
0\longrightarrow\pi^*\Lambda^0(\C)\xrightarrow{\ \alpha\ }\pi^*\Lambda^1(\C)\longrightarrow 0,
\]
其中 $\pi^*\Lambda^0(\C)\iso\pi^*\Lambda^1(\C)\iso\C\times\C$，且 $\alpha:(z,v)\mapsto(z,zv)$。记 $D^-=\{z\mid |z|\le1\}$、$D^+=\{z\mid |z|\ge1\}$。通过包含 $\C\subset\C^+\iso S^2$，可把 $D^-$ 与 $B^-$ 等同，把 $D^+\cup\infty$ 与 $B^+$ 等同。因为 $E_0$ 是 $S^2$ 上的平凡丛，所以 $\pi^*\Lambda^0(\C)=E_0|_\C$。又可构造同构 $\beta:\pi^*\Lambda^1(\C)\to E_{-1}|_\C$，取
\[
\beta:(z,v)\mapsto\begin{cases}(z,v)\in B^-\times\C & z\in D^-,\\(z,z^{-1}v)\in B^+\times\C & z\in D^+.\end{cases}
\]
把圆周 $S^1=\R/2\pi\Z$ 与 $\C$ 中的单位圆等同，注意映射 $f(\theta)=e^{-i\theta}$ 等价于映射 $f(z)=z^{-1}$，故由粘合构造知 $\beta$ 连续。于是外复形 $\Lambda(\C)$ 可改写成 $\C$ 上的复形
\[
0\longrightarrow E_0|_\C\xrightarrow{\ \gamma\ }E_{-1}|_\C\longrightarrow 0,
\]
其中 $\gamma=\beta\circ\alpha$ 由
\[
\gamma:(z,v)\mapsto\begin{cases}(z,zv)\in B^-\times\C & z\in D^-,\\(z,v)\in B^+\times\C & z\in D^+.\end{cases}
\]
给出。最后令 $\gamma(\infty,v)=(\infty,v)$，可把它延拓为 $S^2$ 上的复形
\[
0\longrightarrow E_0\xrightarrow{\ \gamma\ }E_{-1}\longrightarrow 0,
\]
而 $S^2$ 紧，故 $\lambda_1=[E_0]-[E_{-1}]\in\tilde K(S^2)=K(\R^2)$。计算 $\alpha(\lambda_1)=\alpha[E_0]-\alpha[E_{-1}]=0-1=-1$，上面的命题给出 $\aind_{S^1}(\lambda_1)=-\alpha(\lambda_1)=1$。这正是 §3 中完成归一化公理证明所需的结果。
\end{example}

\transnote{符号约定要看清：由构造 $\alpha[E_n]=-n$，而 $\Ind P_n=-n$，两者一致；生成元 $b=[E_{-1}]-[E_0]$ 给出 $\alpha(b)=0-(-1)=1$。后文 Bott 周期性定理的 $\alpha(b)=1$ 与归一化公理都依赖这一符号。若把 $\alpha$ 改成取 $+\deg f$，整条链上的符号会同时反转。}

\subsection{Bott 周期性定理}

取 $b=-\lambda_1=[E_{-1}]-[E_0]$ 作为 $K(\R^2)$ 的生成元，于是 $\alpha(b)=1$。还注意到，由构造有交换图表
\[
\begin{tikzcd}[row sep=large, column sep=large]
K(\R^2\times X)\otimes K(Y) \arrow[r,"\BT"] \arrow[d,"\alpha_x\otimes1"'] & K(\R^2\times X\times Y) \arrow[d,"\alpha_{X\times Y}"] \\
K(X)\otimes K(Y) \arrow[r,"\BT"] & K(X\times Y)
\end{tikzcd}
\]

% ===== 原书 p.44 =====
它说明：对任意 $x\in K(\R^2\times X)$ 与 $y\in K(Y)$，有乘积公式
\[
\alpha(X)\cdot y=\alpha(x\cdot y).
\]
还注意到，映射 $\alpha:K(\R^2\times X)\to K(X)$ 可通过单点紧化 $X^+$ 扩充到局部紧空间 $X$，乘积公式仍成立。现在，Bott 周期性定理可由 K 理论的几条简单代数性质直接推出。

\begin{theorem}[Bott 周期性]
设 $X$ 是局部紧空间，则如下两个同态
\[
\begin{gathered}
\beta:K(X)\to K(\R^2\times X)\\
\alpha:K(\R^2\times X)\to K(X)
\end{gathered}
\]
互为逆。
\end{theorem}

\begin{proof}
先由乘积公式，对任意 $x\in K(X)$ 有
\[
\alpha\beta(X)=\alpha(b\cdot x)=\alpha(b)\cdot x=1\cdot x=x,
\]
故 $\alpha$ 是 $\beta$ 的左逆。

还要证明 $\alpha$ 也是 $\beta$ 的右逆，即对每个 $u\in K(\R^2\times X)$ 验证 $b\cdot\alpha(u)=\beta\alpha(u)=u$。由乘积公式有 $\alpha(u\cdot b)=\alpha(u)\cdot b$，其中 $u\cdot\lambda_1\in K(\R^2\times X\times\R^2)$，我们希望交换这两个乘积。为此考虑映射
\[
\tau:\R^2\times X\times\R^2\to\R^2\times X\times\R^2,
\]
它交换两个 $\R^2$ 的次序。取 $X$ 为一点，相应的变换 $\tau:\R^4\to\R^4$ 由矩阵
\[
\begin{pmatrix}0&I_2\\I_2&0\end{pmatrix},\qquad\text{其中 } I_2=\begin{pmatrix}1&0\\0&1\end{pmatrix}
\]
给出；其行列式为 $+1$，故它与单位元在 $\GL(4,\R)$ 的同一连通分支中。对一般的 $X$，同样可以构造 $\tau$ 与恒等映射之间的同伦，从而 K 理论上的诱导映射
\[
\tau^*:K(\R^2\times X\times\R^2)\to K(\R^2\times X\times\R^2)
\]
是恒等映射。设 $\tilde u\in K(X\times\R^2)$ 是 $u$ 在典范同构 $K(\R^2\times X)\to K(X\times\R^2)$ 下的对应元素，得到
\[
u\cdot b=\tau^*(u\cdot b)=b\cdot\tilde u,
\]
两边作用 $\alpha$，得
\[
\alpha(u)\cdot b=\alpha(u\cdot b)=\alpha(b\cdot\tilde u)=\tilde u\in K(X\times\R^n).
\]
再经由同构 $K(X\times\R^2)\to K(\R^2\times X)$ 变回去，得到
\[
u=b\cdot\alpha(u)=\beta\alpha(u)\in K(\R^2\times X),
\]
证明完成。
\end{proof}

\transnote{关键是 $\det\tau=+1$：两块 $\R^2$ 互换的矩阵行列式为 $\det(-I_2)=+1$，故 $\tau$ 落在 $\GL(4,\R)$ 的单位分支，$\tau^*$ 才是恒等映射，两个乘积才能交换。维数 $2$ 在此处省不掉：若把两个 $\R^2$ 换成两个 $\R^1$，交换矩阵行列式为 $-1$，$\tau^*$ 不再平凡，$\alpha$ 只能是右逆的反号版本。原文先取 $X$ 为一点算行列式，再靠对 $X$ 的自然性回到一般情形。}

% ===== 原书 p.45 =====
\subsection{Thom 同构定理}

这一 Bott 周期性证明因构造是「模块化」的，容易推广。事实上，我们只需要知道逆同态 $\alpha_X:K(\R^2\times X)\to K(X)$ 满足下面两条初等性质：

\textbf{乘积公式}\quad 对一切 $x\in K(\R^2\times X)$ 与 $y\in K(Y)$，$\alpha_{X\times Y}(x\cdot y)=\alpha_X(x)\cdot y$。

\textbf{归一化公理}\quad $\alpha_\pt(b)=1$，其中 $b=\beta_\pt(1)\in K(\R^2)$。

因此，若换一种构造给出逆同态 $\alpha_X:K(\R^2\times X)\to K(X)$，只要它仍满足上述两条性质，我们的证明依然成立。另一种这样的构造见 [7]。

现在把 Bott 周期性的证明推广到等变情形。设给定紧 Lie 群 $G$。取有限维复 $G$-模 $V$，可构造元素 $\lambda_V\in K_G(V)$，它由外复形 $\Lambda(V)$ 给出。把讨论中的 $\R^2$ 换成 $V$，得到同态
\[
\varphi:K_G(X)\xrightarrow{\ \lambda_V^*\BT\cdot\ }K_G(V\times X),
\]
它由 Thom 类 $\lambda_V$ 的对偶丛 $\lambda_V^*$ 作外乘给出。要证这个映射是同构，须构造逆同态 $\alpha_X:K_G(V\times X)\to K_G(X)$，使它满足如下两条类似公理：

\textbf{乘积公式}\quad 对一切 $x\in K_G(V\times X)$ 与 $y\in K_G(Y)$，$\alpha_{X\times Y}(x\cdot y)=\alpha_X(x)\cdot y$。

\textbf{归一化公理}\quad $\alpha_\pt(\lambda_V^*)=1\in R(G)$，其中 $\lambda_V^*=\varphi_\pt(1)\in K_G(V)$。

这样的构造同样可以用一族椭圆算子的指标做到（见 [4, §4]），于是得到

\begin{theorem}[等变周期性]
若 $X$ 是紧 $G$-空间，$V$ 是有限维复 $G$-模，则由 $\lambda_V^*\in K_G(V)$ 外乘给出的同态
\[
\varphi:K_G(X)\to K_G(V\times X)
\]
是同构。
\end{theorem}

本定理有一个特例：取 $G=\U(n)$ 为酉群，$V=\C^n$ 为平凡表示。此时 $X$ 上每个复向量丛都可写成 $E=Y\times_{\U(n)}\C^n$，其中 $Y$ 是自由 $\U(n)$-空间，$Y/\U(n)=X$。回顾 §4：若 $X$ 是自由 $G$-空间，则 $K_G(X)\iso K(X/G)$。于是
\[
K_{\U(n)}(Y)\iso K(X),\qquad K_{\U(n)}(Y\times_{\U(n)}\C^n)\iso K(E).
\]
注意 Thom 类 $\lambda_E\in K(E)$ 对应元素 $\varphi(\lambda_V)\in K_G(V\times X)$，由等变周期性定理立刻得到

% ===== 原书 p.46（上半页，仍属附录 A） =====
\begin{corollary}[Thom 同构定理]
若 $E$ 是紧空间 $X$ 上的复向量丛，则用 Thom 类 $\lambda_E\in K(E)$ 作乘法诱导出同构 $K(X)\to K(E)$。
\end{corollary}

还可以引入第二个群 $H$ 的作用，讨论不受影响。若 $X$ 是 $H$-空间，$E$ 是 $X$ 上的 $H$-丛，则可把 $E$ 写成 $E=Y\times_{\U(n)}\C^n$，其中 $Y$ 是 $H\times\U(n)$-空间，$Y/\U(n)=X$。于是
\[
K_{H\times\U(n)}(Y)\iso K_H(X),\qquad K_{H\times\U(n)}(Y\times_{\U(n)}\C^n)\iso K_H(E),
\]
从而得到等变推广。

\begin{corollary}[等变 Thom 同构定理]
若 $E$ 是紧 $H$-空间 $X$ 上的复 $H$-向量丛，则用 Thom 类 $\lambda_E\in K_H(E)$ 作乘法诱导出同构 $K_H(X)\to K_H(E)$。
\end{corollary}
